\documentclass[11pt]{article}

\usepackage[margin=1in]{geometry}
\usepackage{amsmath,amssymb}
\usepackage{booktabs,tabularx}
\usepackage{enumitem}
\usepackage[dvipsnames]{xcolor}
\usepackage[most]{tcolorbox}
\usepackage{microtype}
\usepackage[hidelinks]{hyperref}

\definecolor{accent}{HTML}{244A73}
\definecolor{paleblue}{HTML}{F1F6FA}
\newcolumntype{Y}{>{\raggedright\arraybackslash}X}
\newtcolorbox{keybox}[1][Key points]{
  colback=paleblue,colframe=accent,fonttitle=\bfseries,
  title=#1,boxrule=0.6pt,arc=1.5mm,left=2mm,right=2mm,top=1mm,bottom=1mm}
\newcommand{\E}{\operatorname{E}}
\newcommand{\Var}{\operatorname{Var}}
\newcommand{\SE}{\operatorname{SE}}

\hypersetup{pdftitle={Ten Essential Concepts in Statistics},
  pdfauthor={Study Notes}}
\setlist{nosep}
\setcounter{tocdepth}{2}

\title{\textbf{Ten Essential Concepts in Statistics}\\[0.4em]
  \large A concise guide for introductory students}
\author{}
\date{}

\begin{document}
\maketitle

\begin{abstract}
Statistics provides a language for learning from data while acknowledging
uncertainty.  These notes develop ten foundational ideas in a logical
sequence: classifying data, describing it numerically and graphically,
modeling chance, making inferences, and studying relationships.  Definitions,
worked examples, assumptions, and interpretations are emphasized throughout.
\end{abstract}

\tableofcontents
\newpage

\section{Types of Data}\label{sec:data}
A \emph{variable} is a characteristic recorded for each observational unit;
its recorded values form data.  The kind of variable determines which
summaries, graphs, and statistical procedures are meaningful.  Before
classifying variables, identify where the data came from and what population
the data are meant to represent.

\subsection{Data structure and sources}
A \emph{population} is the complete group or process about which conclusions
are desired.  A \emph{sample} is the subset actually observed.  An
\emph{observational unit} is the person, object, time period, or case on
which measurements are recorded.  A \emph{data value} is one recorded value
of one variable for one unit.  A rectangular \emph{data matrix} usually puts
observational units in rows and variables in columns.  A \emph{sampling
frame} is the list or mechanism from which the sample is selected; if the
frame omits relevant units, generalizability is limited.

Table~\ref{tab:data-matrix} is a hypothetical classroom data matrix.  If the
research question concerns all students in a university statistics course,
the population might be all enrolled students, and a possible sampling frame
is the registrar's course roster.

\begin{table}[ht]
\centering
\caption{A hypothetical five-student data matrix.}\label{tab:data-matrix}
\begin{tabularx}{\textwidth}{@{}lYYYYY@{}}
\toprule
Student & Major & Academic year & Study hours & Exam score & Attended review? \\
\midrule
A & Economics & First & 2 & 58 & No \\
B & Biology & Second & 4 & 71 & Yes \\
C & History & First & 3 & 65 & No \\
D & Mathematics & Third & 6 & 82 & Yes \\
E & Economics & Second & 5 & 76 & Yes \\
\bottomrule
\end{tabularx}
\end{table}

In this table, major and review attendance are nominal categorical variables,
academic year is ordinal, and study hours and exam score are quantitative.
Study hours are continuous in an ideal measurement model but may be rounded
when recorded; exam score is quantitative on the test scale.

\subsection{From data to conclusions}
\emph{Simple random sampling} gives each possible sample of the desired size
the same chance of selection.  \emph{Stratified sampling} first divides the
population into important groups, such as academic years, and then samples
within each group.  \emph{Cluster sampling} randomly selects natural groups,
such as sections or dormitories, and observes units within selected groups.
\emph{Convenience sampling} uses units that are easy to reach; it is often
fast but can be badly biased.

An \emph{observational study} records variables without assigning treatments.
A \emph{randomized experiment} assigns treatments by a chance mechanism.
Random sampling supports generalization to a suitably represented population;
random assignment supports causal conclusions under appropriate experimental
controls.  Neither mechanism automatically guarantees the other.

\begin{table}[ht]
\centering
\caption{Design features and the conclusions they support.}\label{tab:design}
\begin{tabularx}{\textwidth}{@{}lYY@{}}
\toprule
Design feature & Main support & Limitation \\
\midrule
Random sample from a good frame & Generalizing sample patterns to the frame population & Does not by itself prove cause and effect \\
Random assignment in an experiment & Comparing treatments causally, if controls are sound & Does not by itself ensure the subjects represent a wider population \\
Convenience or volunteer sample & Quick description of responding units & Vulnerable to selection bias and weak generalizability \\
Observational comparison & Association in naturally occurring data & Confounding and reverse direction can remain \\
\bottomrule
\end{tabularx}
\end{table}

\emph{Selection bias} occurs when the way units enter the sample systematically
favors some outcomes.  \emph{Nonresponse bias} occurs when selected units who
do not respond differ from responders.  \emph{Measurement bias} occurs when
the measurement process systematically overstates or understates values.
\emph{Confounding} occurs when an outside variable is related to both an
explanatory variable and a response, making a causal explanation ambiguous.
For example, a volunteer survey about sleep in an early-morning class may
overrepresent students willing to attend early classes and underrepresent
absent or sleep-deprived students.  A random sample from the registrar list
would usually support broader generalization, although nonresponse and
measurement quality still need attention.

\subsection{Qualitative and quantitative variables}
\emph{Qualitative} (categorical) data place units into groups.  Their labels
need not be numbers.  \emph{Quantitative} data record numerical amounts or
counts for which arithmetic has a meaningful interpretation.

\begin{table}[ht]
\centering
\caption{The principal data types.}\label{tab:data-types}
\begin{tabularx}{\textwidth}{@{}lYY@{}}
\toprule
Type & Definition & Examples \\
\midrule
Qualitative & Category or quality rather than a numerical magnitude &
Favorite sport; a movie rating; major code 01 \\
Quantitative & Numerical count or measurement &
Number of goals; height; temperature; exam score \\
\bottomrule
\end{tabularx}
\end{table}

\subsubsection{Nominal and ordinal categories}
Nominal categories have no inherent ranking.  For example, favorite sport
might be hockey, golf, or soccer; no category is intrinsically above another.
Ordinal categories have an order, but gaps between categories are not
necessarily equal.  Movie ratings \emph{bad}, \emph{fair}, and \emph{good}
are ordered, yet the distance from bad to fair is not a measurable numerical
interval.  Numeric codes can still represent categorical labels: jersey
numbers, identification numbers, and major codes are not quantitative
measurements.  Arithmetic with ordinal codes, such as averaging satisfaction
ratings coded 1 through 5, requires caution unless an interval-scale
interpretation is defensible.

\begin{table}[ht]
\centering
\caption{Two forms of qualitative data.}
\begin{tabularx}{\textwidth}{@{}lYY@{}}
\toprule
Scale & Ordering? & Appropriate comparisons \\
\midrule
Nominal & No natural order & Equal/not equal; category frequencies \\
Ordinal & Natural order & Equal/not equal and before/after; ranks \\
\bottomrule
\end{tabularx}
\end{table}

\subsubsection{Discrete and continuous quantities}
A discrete variable has a finite or countably infinite set of possible
values, often arising from counting.  Numbers of goals and students are
discrete: fractional students are not possible.  A continuous variable can,
in an ideal measurement model, take values over a finite interval, a
semi-infinite interval, the whole real line, or a union of intervals.  Height
and temperature are continuous even when an instrument rounds their recorded
values, because continuous measurement models are idealizations of the
underlying quantity.

\begin{table}[ht]
\centering
\caption{Two forms of quantitative data.}
\begin{tabularx}{\textwidth}{@{}lYY@{}}
\toprule
Type & Possible values & Typical origin \\
\midrule
Discrete & Finite or countably infinite set & Counting, such as goals \\
Continuous & Values in intervals under a measurement model & Measuring, such as height \\
\bottomrule
\end{tabularx}
\end{table}

\begin{keybox}
Data type guides numerical summaries, graphs, and procedures.  Study design
guides what conclusions are justified: random sampling concerns
representativeness, while random assignment concerns causation.
\end{keybox}

\section{Measures of Central Tendency}\label{sec:center}
A \emph{measure of central tendency} is a numerical description of a typical
or central value in a distribution.  The mean, median, and mode capture
different notions of center.

\subsection{Mean}
For observations $x_1,x_2,\ldots,x_n$, the sample mean is
\begin{equation}\label{eq:mean}
 \bar{x}=\frac{1}{n}\sum_{i=1}^{n}x_i.
\end{equation}
Here $x_i$ is observation $i$, $n$ is the sample size, $\sum$ denotes
addition over all observations, and $\bar{x}$ is the sample mean.  The
population analogue is
\[
\mu=\frac{1}{N}\sum_{i=1}^{N}x_i,
\]
where $N$ is population size and $\mu$ is the population mean.

For seven illustrative test scores
\[
72,\ 78,\ 80,\ 85,\ 85,\ 86,\ 86,
\]
the sum is $572$, so
\[
 \bar{x}=\frac{572}{7}=81.714\ldots\approx81.7\text{ points}.
\]
The mean is a balance point and uses every value, so an extreme score can
move it substantially.

\subsection{Weighted and frequency means}
When observations have weights $w_i$, the weighted mean is
\begin{equation}\label{eq:weighted-mean}
\bar{x}_w=\frac{\sum_i w_i x_i}{\sum_i w_i}.
\end{equation}
For a frequency table with value $x_i$ occurring $f_i$ times,
\begin{equation}\label{eq:freq-mean}
\bar{x}=\frac{\sum_i f_i x_i}{\sum_i f_i}.
\end{equation}
For quiz scores $6,7,8,9$ with frequencies $1,2,3,2$, respectively,
\[
\bar{x}=\frac{1(6)+2(7)+3(8)+2(9)}{1+2+3+2}
       =\frac{62}{8}=7.75.
\]

\begin{table}[ht]
\centering
\caption{A reproducible frequency table for a mean.}\label{tab:freq-mean}
\begin{tabularx}{0.7\textwidth}{@{}lYYY@{}}
\toprule
Score $x_i$ & Frequency $f_i$ & Product $f_ix_i$ & Interpretation \\
\midrule
6 & 1 & 6 & one score of 6 \\
7 & 2 & 14 & two scores of 7 \\
8 & 3 & 24 & three scores of 8 \\
9 & 2 & 18 & two scores of 9 \\
\midrule
Total & 8 & 62 & $\bar{x}=62/8$ \\
\bottomrule
\end{tabularx}
\end{table}

\subsection{Median and mode}
After sorting $n$ observations, the \emph{median} is the middle observation
when $n$ is odd and the mean of the two middle observations when $n$ is even.
For the ordered odd-size data
\[
62,\ 70,\ 70,\ 85,\ 90,
\]
the median is the third value, $70$.  For the ordered even-size data
\[
62,\ 70,\ 85,\ 90,
\]
the median is $(70+85)/2=77.5$.

The \emph{mode} is the most frequent value or category.  The data
$62,70,70,85,90$ have mode $70$.  The data $1,2,3,4$ have no mode if a mode
is required to occur more often than the rest.  The data $1,1,2,3,3,4$ are
bimodal, with modes $1$ and $3$.

\begin{table}[ht]
\centering
\caption{Choosing a measure of center.}\label{tab:center}
\begin{tabularx}{\textwidth}{@{}lYYY@{}}
\toprule
Measure & Definition & Strength & Caution \\
\midrule
Mean & Arithmetic average & Uses every value & Sensitive to outliers and skew \\
Median & Middle ordered value & Resistant to outliers & Ignores magnitudes away from center \\
Mode & Most frequent value & Works for categories & May be absent or nonunique \\
\bottomrule
\end{tabularx}
\end{table}

\begin{keybox}[Common cautions]
Always report units.  A mean need not be an observed value.  For a strongly
skewed quantitative distribution, the median often better represents a
typical observation.  Right-skewed distributions often have mean greater than
median, and left-skewed distributions often have mean less than median, but
this is a tendency rather than a universal theorem.  For nominal data, the
mode is the only one of these three measures that applies.  Center alone is
incomplete, motivating spread.
\end{keybox}

\section{Measures of Spread}\label{sec:spread}
A \emph{measure of spread} describes how far observations lie from one
another or from their center.  Two data sets can share a mean yet differ
greatly in variability.

\subsection{Range, percentiles, quartiles, and IQR}
For ordered data, the range and interquartile range (IQR) are
\begin{align}
 \text{range}&=x_{\max}-x_{\min},\label{eq:range}\\
 \mathrm{IQR}&=Q_3-Q_1,\label{eq:iqr}
\end{align}
where $x_{\max}$ and $x_{\min}$ are the maximum and minimum, $Q_1$ is the
first quartile, $Q_2$ is the median, and $Q_3$ is the third quartile.  A
percentile is a location in an ordered distribution; roughly, the $p$th
percentile has about $p\%$ of observations at or below it.  The
\emph{five-number summary} is
\[
\text{minimum},\quad Q_1,\quad \text{median},\quad Q_3,\quad \text{maximum}.
\]
These notes use the median-of-halves convention: after finding the median,
$Q_1$ is the median of the lower half and $Q_3$ is the median of the upper
half, excluding the overall median when $n$ is odd.  Software may use other
legitimate percentile conventions, so the convention should be stated.

For ordered study hours
\[
1,\ 2,\ 4,\ 5,\ 5,\ 6,\ 7,\ 10,
\]
the sum is $40$, so $\bar{x}=40/8=5$ hours.  The range is $10-1=9$ hours.
Using the median-of-halves convention,
\[
 Q_1=\frac{2+4}{2}=3,\qquad
 Q_2=\frac{5+5}{2}=5,\qquad
 Q_3=\frac{6+7}{2}=6.5,
\]
and therefore $\mathrm{IQR}=6.5-3=3.5$ hours.  The middle half of the
observations spans $3.5$ hours.

\subsection{Outlier fences and box-plot conventions}
A common rule flags observations below or above the fences
\begin{equation}\label{eq:iqr-fences}
Q_1-1.5\,\mathrm{IQR},\qquad Q_3+1.5\,\mathrm{IQR}.
\end{equation}
For the study-hour data, the fences are
\[
3-1.5(3.5)=-2.25,\qquad 6.5+1.5(3.5)=11.75.
\]
No observation is outside these fences.  A min--max box plot would use
whiskers at $1$ and $10$.  A Tukey-style box plot also uses whiskers at the
most extreme non-outlying observations; here those are again $1$ and $10$.
If a value such as $18$ hours were added, the quartiles and fences would need
to be recomputed; outlying values would be plotted separately rather than
used as Tukey whisker endpoints.

\subsection{Variance, standard deviation, and standardized values}
Population variance and standard deviation are
\begin{equation}\label{eq:pop-sd}
 \sigma^2=\frac{1}{N}\sum_{i=1}^{N}(x_i-\mu)^2,
 \qquad \sigma=\sqrt{\sigma^2}.
\end{equation}
Sample variance and standard deviation are
\begin{equation}\label{eq:samp-sd}
 s^2=\frac{1}{n-1}\sum_{i=1}^{n}(x_i-\bar{x})^2,
 \qquad s=\sqrt{s^2}.
\end{equation}
Dividing by $n-1$ makes $s^2$ an unbiased estimator of population variance
under random sampling.  Once $\bar{x}$ is estimated, the deviations sum to
zero, leaving $n-1$ independently variable deviations.

For the study-hour sample, the squared deviations from $5$ are
\[
16,\ 9,\ 1,\ 0,\ 0,\ 1,\ 4,\ 25,
\]
whose sum is $56$.  Hence
\[
s^2=\frac{56}{8-1}=8,
\qquad s=\sqrt{8}=2.828\ldots\approx2.8\text{ hours}.
\]
Standard deviation is a root-mean-square scale of deviations from the mean,
not the mean absolute distance from the mean.  There is no universal rule
that a fixed percentage of arbitrary data lies within one standard deviation;
about $68\%$ within one standard deviation is a normal-model result.

A descriptive sample standardized value is
\begin{equation}\label{eq:sample-z}
z_i=\frac{x_i-\bar{x}}{s}.
\end{equation}
For the $10$-hour observation, $z=(10-5)/2.828\approx1.77$, so it is about
$1.77$ sample standard deviations above the sample mean.

\begin{table}[ht]
\centering
\caption{Comparison of spread measures.}
\begin{tabularx}{\textwidth}{@{}lYYY@{}}
\toprule
Measure & Formula & Uses & Sensitivity \\
\midrule
Range & $x_{\max}-x_{\min}$ & Full span & Very sensitive to extremes \\
IQR & $Q_3-Q_1$ & Middle $50\%$ & Resistant to extremes \\
Standard deviation & Root mean squared scale of deviations & Spread about mean & Sensitive to extremes \\
\bottomrule
\end{tabularx}
\end{table}

\begin{keybox}
Variance has squared units; standard deviation has the original units.  Pair
mean with standard deviation for roughly symmetric data and median with IQR
for skewed data.  Graphs in Section~\ref{sec:graphs} reveal shape that no
single spread statistic can show.
\end{keybox}

\section{Graphical Representations of Data}\label{sec:graphs}
A \emph{statistical graph} maps data values or categories to visual marks so
that frequencies, proportions, distributional shape, and unusual values can
be seen.

\subsection{Categorical displays and frequency tables}
For favorite subject, the frequencies are Mathematics $12$, Science $9$,
English $3$, and History $6$, for a total of $30$.  A frequency table gives
counts, while a relative-frequency table divides each count by the total.

\begin{table}[ht]
\centering
\caption{Favorite-subject frequency and relative-frequency table.}\label{tab:fav-subject}
\begin{tabularx}{0.75\textwidth}{@{}lYY@{}}
\toprule
Subject & Frequency & Relative frequency \\
\midrule
Mathematics & 12 & $12/30=0.40=40\%$ \\
Science & 9 & $9/30=0.30=30\%$ \\
English & 3 & $3/30=0.10=10\%$ \\
History & 6 & $6/30=0.20=20\%$ \\
\midrule
Total & 30 & $1.00=100\%$ \\
\bottomrule
\end{tabularx}
\end{table}

A bar chart gives each category a separated bar with height equal to its
frequency or relative frequency.  A pie chart uses sector proportions; for
example Mathematics occupies $40\%$ of the circle.  Cumulative relative
frequency is meaningful for ordered categories or quantitative intervals, but
not for unordered subjects unless an order is imposed for a special purpose.

\subsection{Quantitative displays and distribution shape}
A histogram divides a quantitative scale into intervals (bins) and uses
touching bars to display frequencies or densities.  Bin widths and boundaries
must be reported because they influence appearance.  A dotplot places one dot
for each observation along a number line and is useful for small samples; a
stem-and-leaf plot is another compact hand display.  A box plot summarizes
the five-number summary and is especially effective for comparing groups.

\begin{keybox}[Describe a distribution]
For a quantitative distribution, describe shape, center, spread, and unusual
features in context.  Shape includes symmetry, right or left skew, unimodality
or bimodality, clusters, gaps, and outliers or contextual anomalies.
\end{keybox}

Compact dotplot sketches can show the main shape ideas without extra
packages:
\[
\begin{array}{c|l}
\text{Roughly symmetric} & 1{:}\bullet\quad 2{:}\bullet\bullet\quad 3{:}\bullet\bullet\bullet\quad 4{:}\bullet\bullet\quad 5{:}\bullet \\
\text{Right-skewed} & 1{:}\bullet\bullet\bullet\quad 2{:}\bullet\bullet\quad 3{:}\bullet\quad 6{:}\bullet \\
\text{Bimodal} & 1{:}\bullet\bullet\quad 2{:}\bullet\quad 5{:}\bullet\quad 6{:}\bullet\bullet
\end{array}
\]
The first sketch is balanced around its center, the second has a long right
tail, and the third has two clusters separated by a gap.

For student heights, the supplied five-number summary is
\[
151,\quad158.5,\quad163.5,\quad169,\quad178.
\]
Thus the box extends from $Q_1=158.5$ to $Q_3=169$, a line marks the median
$163.5$, and whiskers reach the minimum $151$ and maximum $178$ under the
simple min--max convention.  Under a Tukey convention, outlier fences require
raw data or enough information to identify non-outlying extremes.

\begin{table}[ht]
\centering
\caption{Common statistical displays.}\label{tab:graphs}
\begin{tabularx}{\textwidth}{@{}lYYY@{}}
\toprule
Display & Data type & Main purpose & Important caution \\
\midrule
Bar chart & Categorical & Compare counts/proportions & Bars are separated; use a common baseline \\
Pie chart & Categorical & Show parts of one whole & Poor for many or similar-sized categories \\
Dotplot & Quantitative & Show individual values in small samples & Can become crowded for large samples \\
Histogram & Quantitative & Show shape, center, and spread & Conclusions depend on bins \\
Box plot & Quantitative & Show quartiles and compare groups & Conceals modes and detailed shape \\
Scatterplot & Paired quantitative & Show form, direction, strength, clusters, and outliers & Association need not be linear or causal \\
\bottomrule
\end{tabularx}
\end{table}

\begin{keybox}[Common cautions]
Do not use a bar chart as though categories were a continuous numerical
scale, or a histogram for category labels.  Axes, units, binning, and the
sample represented must be clear.  Visualized relative frequencies lead
naturally to probability.
\end{keybox}

\section{Probability Fundamentals}\label{sec:probability}
Probability is a number from $0$ to $1$ that quantifies uncertainty about an
event: $0$ means impossible and $1$ means certain.  Equivalent forms include
\[
0.7=70\%=\frac{7}{10}.
\]

\subsection{Experiments, outcomes, events, and axioms}
A random experiment has an uncertain result.  Its \emph{sample space} $S$ is
the set of possible outcomes, and an \emph{event} $A$ is a subset of $S$.
The basic probability axioms are
\begin{equation}\label{eq:prob-axioms}
0\le P(A)\le1,\qquad P(S)=1,
\end{equation}
and if $A$ and $B$ are disjoint events, then
\[
P(A\cup B)=P(A)+P(B).
\]
When all elementary outcomes are equally likely,
\begin{equation}\label{eq:classical-prob}
 P(A)=\frac{\text{number of outcomes favorable to }A}
 {\text{total number of outcomes in }S}.
\end{equation}
This favorable-over-total formula is not valid for unequal elementary
outcomes unless the outcomes have first been weighted appropriately.  For a
fair six-sided die, $P(\text{three})=1/6$.

The complement and general addition rules are
\begin{align}
P(A^c)&=1-P(A),\label{eq:complement}\\
P(A\cup B)&=P(A)+P(B)-P(A\cap B),\label{eq:addition}
\end{align}
where $A^c$ means ``not $A$,'' $A\cup B$ means ``$A$ or $B$ (or both),''
and $A\cap B$ means ``both.''

\subsection{Conditional probability, multiplication, and independence}
Conditional probability is
\begin{equation}\label{eq:conditional}
P(A\mid B)=\frac{P(A\cap B)}{P(B)},\qquad P(B)>0.
\end{equation}
Rearranging gives the general multiplication rule
\begin{equation}\label{eq:multiplication}
P(A\cap B)=P(A\mid B)P(B)=P(B\mid A)P(A).
\end{equation}
Events $A$ and $B$ are independent if knowing one occurred does not change
the probability of the other; equivalently,
\begin{equation}\label{eq:independence}
P(A\cap B)=P(A)P(B).
\end{equation}
Independent fair die rolls therefore give
\[
P(\text{two consecutive threes})=\frac16\cdot\frac16=\frac1{36}.
\]

Mutually exclusive events cannot occur together; independent events do not
influence one another.  They are not synonymous.  On one die roll, $A=$
``roll a 1'' and $B=$ ``roll a 2'' are mutually exclusive, but because
$P(A\cap B)=0\ne(1/6)(1/6)$, they are not independent.  On two die rolls,
$A=$ ``first roll is 1'' and $B=$ ``second roll is 1'' are independent but
not mutually exclusive, because both can occur.

\subsection{Total probability and Bayes' theorem}
If $B$ and $B^c$ split the sample space, then
\begin{equation}\label{eq:total-prob}
P(A)=P(A\mid B)P(B)+P(A\mid B^c)P(B^c).
\end{equation}
Bayes' theorem reverses a conditional probability:
\begin{equation}\label{eq:bayes}
P(B\mid A)=\frac{P(A\mid B)P(B)}{P(A)}.
\end{equation}

Suppose a rare condition affects $1\%$ of a population.  A screening test is
positive for $90\%$ of people with the condition and for $5\%$ of people
without it.  For a hypothetical group of $10{,}000$ people, the expected
counts are:

\begin{table}[ht]
\centering
\caption{A $2\times2$ table illustrating Bayes' theorem and base rates.}\label{tab:bayes}
\begin{tabularx}{0.85\textwidth}{@{}lYYY@{}}
\toprule
Group & Positive test & Negative test & Total \\
\midrule
Condition & 90 & 10 & 100 \\
No condition & 495 & 9405 & 9900 \\
\midrule
Total & 585 & 9415 & 10000 \\
\bottomrule
\end{tabularx}
\end{table}

Thus $P(\text{positive})=585/10000=0.0585$.  The reversed conditional
probability is
\[
P(\text{condition}\mid\text{positive})=\frac{90}{585}=0.1538\ldots\approx15.4\%.
\]
Even with a sensitive test, most positive tests are false positives because
the condition is rare; this is the base-rate effect.

\subsection{Theoretical and experimental probability}
Experimental probability, or relative frequency, estimates probability from
observed repetitions:
\begin{equation}\label{eq:empirical-prob}
 \widehat{P}(A)=\frac{\text{number of observed occurrences of }A}
 {\text{number of trials}}.
\end{equation}
If two threes occur in $20$ rolls, $\widehat P(\text{three})=2/20=0.10$.
This differs from $1/6\approx0.1667$ because short-run samples fluctuate; with
many independent trials under stable conditions, relative frequency tends to
approach the true probability.

\subsection{Counting methods used in probability}
\begin{keybox}[Counting methods]
For a positive integer $n$, $n!=n(n-1)\cdots2\cdot1$ and $0!=1$.
Permutations count ordered arrangements: $P(n,r)=n!/(n-r)!$.  Combinations
count unordered selections: $\binom nr=n!/[r!(n-r)!]$.  The binomial formula
uses combinations because it counts which $k$ of the $n$ trials are successes,
not the order in which the selected positions are named.
\end{keybox}

\begin{keybox}
Probability rules become distribution rules when outcomes are assigned
numerical values.  Discrete distributions use probability masses; continuous
distributions use density and area.
\end{keybox}

\section{Discrete Probability Distributions}\label{sec:discrete}
A \emph{discrete random variable} $X$ maps outcomes to a finite or countably
infinite numerical set.  Its probability mass function (PMF) is
$p(x)=P(X=x)$, with $p(x)\ge0$ and $\sum_xp(x)=1$.  Its cumulative
distribution function (CDF) is
\begin{equation}\label{eq:disc-cdf}
F(x)=P(X\le x)=\sum_{t\le x}p(t).
\end{equation}

\subsection{Expected value, variance, and rules}
For possible values $x$,
\begin{align}
\mu_X=\E(X)&=\sum_x x p(x),\label{eq:disc-mean}\\
\Var(X)&=\sum_x(x-\mu_X)^2p(x),\qquad
\sigma_X=\sqrt{\Var(X)}.\label{eq:disc-var}
\end{align}
Useful rules include
\begin{align}
\E(a+bX)&=a+b\E(X),\label{eq:expect-transform}\\
\Var(a+bX)&=b^2\Var(X),\label{eq:var-transform}
\end{align}
and, when $X$ and $Y$ are independent,
\begin{equation}\label{eq:var-sum}
\Var(X+Y)=\Var(X)+\Var(Y).
\end{equation}

\subsection{Bernoulli and binomial distributions}
A Bernoulli random variable records one trial with success probability $p$:
$X=1$ for success and $X=0$ for failure.  A binomial experiment repeats this
structure a fixed number $n$ of independent trials, each with the same
success probability $p$.  If $X$ counts successes, then
\begin{equation}\label{eq:binomial}
X\sim\operatorname{Binomial}(n,p),\qquad
P(X=k)=\binom nk p^k(1-p)^{n-k},\qquad k=0,\ldots,n.
\end{equation}
Also $\E(X)=np$ and $\Var(X)=np(1-p)$.

Let success mean rolling a six, in three independent fair die rolls.  Then
\[
X\sim\operatorname{Binomial}\left(3,\frac16\right).
\]
Applying \eqref{eq:binomial} produces the complete distribution and CDF:

\begin{table}[ht]
\centering
\caption{PMF and CDF for $X\sim\operatorname{Binomial}(3,1/6)$.}\label{tab:binomial-cdf}
\begin{tabularx}{0.75\textwidth}{@{}lYY@{}}
\toprule
$k$ & $P(X=k)$ & $F(k)=P(X\le k)$ \\
\midrule
0 & $125/216$ & $125/216$ \\
1 & $75/216$ & $200/216$ \\
2 & $15/216$ & $215/216$ \\
3 & $1/216$ & $216/216=1$ \\
\bottomrule
\end{tabularx}
\end{table}

The numerators in the PMF sum to $216$, verifying total probability $1$.
The expected number of sixes is $np=3(1/6)=0.5$; this long-run average need
not be an attainable count in one experiment.  A tail probability can be
computed from the complement:
\[
P(X\ge1)=1-P(X=0)=1-\frac{125}{216}=\frac{91}{216}\approx0.4213.
\]

\subsection{The Poisson model for counts}
A Poisson random variable models counts in a specified interval of time,
space, area, or volume under a constant-rate model.  If the expected count in
the interval is $\lambda>0$, then
\begin{equation}\label{eq:poisson}
P(X=k)=e^{-\lambda}\frac{\lambda^k}{k!},\qquad k=0,1,2,\ldots,
\end{equation}
where $\lambda=\E(X)=\Var(X)$.  Under the model, counts in disjoint intervals
are independent and the expected count is proportional to interval length.
Dependence, clustering, bounded counts, or changing rates can invalidate the
model.

If calls arrive at an average rate of $3$ per hour, then for one hour
$X\sim\operatorname{Poisson}(3)$ and
\[
P(X=0)=e^{-3}\approx0.0498,
\qquad
P(X\ge1)=1-e^{-3}\approx0.9502.
\]
For a half-hour interval under the same constant rate, the expected count is
$\lambda=1.5$ rather than $3$.

\begin{keybox}
A valid PMF assigns probabilities to individual values, and all masses sum
to $1$.  Check binomial and Poisson conditions before using their formulas;
sampling without replacement can violate independence, and changing rates can
invalidate a Poisson model.
\end{keybox}

\section{Continuous Probability Distributions}\label{sec:continuous}
A \emph{continuous random variable} is modeled with probabilities over
intervals.  Its probability density function (PDF) $f$ satisfies $f(x)\ge0$
and
\begin{equation}\label{eq:density}
P(a\le X\le b)=\int_a^b f(x)\,dx,
\qquad \int_{-\infty}^{\infty}f(x)\,dx=1.
\end{equation}
Height is not itself probability: probability is area.  Consequently
$P(X=x)=0$ for every individual value $x$.  The continuous CDF is
\begin{equation}\label{eq:cont-cdf}
F(x)=P(X\le x)=\int_{-\infty}^{x}f(t)\,dt,
\end{equation}
which converts density into cumulative probability.

When the integrals exist,
\begin{align}
\E(X)&=\int_{-\infty}^{\infty}x f(x)\,dx,\label{eq:cont-mean}\\
\Var(X)&=\int_{-\infty}^{\infty}(x-\mu)^2 f(x)\,dx.\label{eq:cont-var}
\end{align}

\begin{table}[ht]
\centering
\caption{Probability mass versus density.}
\begin{tabularx}{\textwidth}{@{}lYY@{}}
\toprule
Feature & PMF $p(x)$ & PDF $f(x)$ \\
\midrule
Variable & Discrete & Continuous \\
Point value & $p(x)=P(X=x)$ may be positive & $P(X=x)=0$ \\
Interval probability & Sum of masses & Area $\int_a^b f(x)\,dx$ \\
Normalization & $\sum_xp(x)=1$ & $\int_{-\infty}^{\infty}f(x)\,dx=1$ \\
\bottomrule
\end{tabularx}
\end{table}

\subsection{Uniform and normal distributions}
Not every continuous distribution is normal.  For example, if a bus arrives
uniformly during a $10$-minute window, then $X\sim\operatorname{Uniform}(0,10)$
has density $f(x)=1/10$ for $0\le x\le10$ and $0$ otherwise.  Therefore
\[
P(2\le X\le5)=\int_2^5\frac{1}{10}\,dx=0.3.
\]

A normal random variable with mean $\mu$ and standard deviation $\sigma>0$
has density
\begin{equation}\label{eq:normal-density}
f(x)=\frac{1}{\sigma\sqrt{2\pi}}
\exp\!\left[-\frac{(x-\mu)^2}{2\sigma^2}\right].
\end{equation}
The bell-shaped curve is symmetric about $\mu$.  Approximately $68\%$,
$95\%$, and $99.7\%$ of its area lies within one, two, and three standard
deviations of $\mu$, respectively.  These empirical-rule percentages are
normal-model approximations; they are not, by themselves, a test of normality.

\subsection{Standardization and inverse normal calculations}
The $z$-score of a value $x$ in a normal population is
\begin{equation}\label{eq:zscore}
z=\frac{x-\mu}{\sigma},
\end{equation}
the signed number of population standard deviations between $x$ and $\mu$.
Suppose adult height $X$ is normally distributed with
\[
\mu=175\text{ cm},\qquad \sigma=6\text{ cm},\qquad x=184\text{ cm}.
\]
Then
\[
z=\frac{184-175}{6}=1.5.
\]
If $Z$ is standard normal and $\Phi(z)=P(Z\le z)$, a standard normal table or
software gives
\[
P(X\le184)=P(Z\le1.5)=\Phi(1.5)\approx0.9332.
\]
Thus about $93.32\%$ of this modeled population is at most $184$ cm tall.

The inverse problem starts with a probability.  The 90th percentile has
$z_{0.90}\approx1.282$, so
\[
x=\mu+z_{0.90}\sigma=175+1.282(6)=182.692\text{ cm}\approx182.7\text{ cm}.
\]
Similarly, $z_{0.975}=1.96$ leaves $2.5\%$ in the upper tail and $2.5\%$ in
the lower tail, giving the central $95\%$ area used in many confidence
intervals.

\begin{keybox}[Common cautions]
A $z$-score is a standardized location, not a probability.  Normality should
not be assumed solely because a variable is continuous.  Areas under sampling
distributions support the confidence intervals of the next section.
\end{keybox}

\section{Confidence Intervals}\label{sec:ci}
A \emph{confidence interval} is a range computed from sample data by a method
designed to capture an unknown population parameter at a stated long-run
rate.  It combines a point estimate with a margin of error.

\subsection{Sampling distributions: the bridge to inference}
A \emph{parameter} is a fixed population quantity, such as $\mu$ or $p$.  A
\emph{statistic} is a sample quantity that varies from sample to sample.  An
\emph{estimator} is the rule or random quantity used to estimate a parameter,
such as $\bar X$ for $\mu$; an \emph{estimate} is the realized value, such as
$\bar x=9.2$.

A \emph{sampling distribution} is the probability distribution of a statistic
computed over repeated samples of the same size from the same population or
process.  If observations are independent with population mean $\mu$ and
standard deviation $\sigma$, then for repeated samples of size $n$,
\begin{equation}\label{eq:xbar-sampling}
\E(\bar X)=\mu,
\qquad
\SE(\bar X)=\frac{\sigma}{\sqrt n}.
\end{equation}
The law of large numbers says sample averages stabilize near the population
mean as sample size increases.  The central limit theorem says that, under
appropriate conditions, the standardized sampling distribution of the sample
mean becomes approximately normal as $n$ grows; it does not say the population
itself becomes normal.

\begin{table}[ht]
\centering
\caption{Standard deviation versus standard error.}\label{tab:sd-se}
\begin{tabularx}{\textwidth}{@{}lYY@{}}
\toprule
Quantity & Describes & Changes with larger $n$? \\
\midrule
Standard deviation & Variability among individual observations & Not necessarily \\
Standard error & Sample-to-sample variability of a statistic & Usually decreases, often like $1/\sqrt n$ \\
\bottomrule
\end{tabularx}
\end{table}

A compact schematic is
\[
\text{population distribution}\longrightarrow
\text{one random sample}\longrightarrow
\bar x\longrightarrow
\text{many repeated }\bar X\text{'s}\longrightarrow
\text{sampling distribution}.
\]
This bridge is why study design from Section~\ref{sec:data} matters for
inference: biased or dependent samples can make the sampling distribution
used by a procedure inappropriate.

\subsection{Statistics and parameters}
\begin{table}[ht]
\centering
\caption{Population quantities and sample summaries.}
\begin{tabularx}{\textwidth}{@{}lYY@{}}
\toprule
Concept & Population parameter & Sample statistic \\
\midrule
Mean & $\mu$, fixed but usually unknown & $\bar{x}$, varies between samples \\
Standard deviation & $\sigma$ & $s$ \\
Proportion & $p$ & $\widehat p$ \\
\bottomrule
\end{tabularx}
\end{table}
A statistic estimates its corresponding parameter.  Its sampling variability
is summarized by a standard error.

\subsection{A general confidence-interval workflow}
A useful workflow is: state the research question, identify the population
parameter, examine sampling design and assumptions, select the procedure,
calculate the estimate and uncertainty, interpret in context, and state
design or model limitations.  Confidence intervals estimate population
parameters or effects while communicating uncertainty and precision; some,
but not all, target parameters are effect-size measures.

\begin{table}[ht]
\centering
\caption{Common standard normal critical values.}\label{tab:z-critical}
\begin{tabularx}{0.55\textwidth}{@{}lY@{}}
\toprule
Confidence level & $z^*$ \\
\midrule
$90\%$ & $1.645$ \\
$95\%$ & $1.960$ \\
$99\%$ & $2.576$ \\
\bottomrule
\end{tabularx}
\end{table}
A central confidence area leaves two equal tails; for $95\%$ confidence,
$2.5\%$ is left in each tail, so $z^*=z_{0.975}=1.96$.

\subsection{A one-sample \texorpdfstring{$z$}{z}-interval for a mean}
When population standard deviation $\sigma$ is known, the interval is
\begin{equation}\label{eq:zinterval}
\bar{x}\pm z^*\frac{\sigma}{\sqrt n}
=\bar{x}\pm \mathrm{ME},
\qquad \SE(\bar{x})=\frac{\sigma}{\sqrt n}.
\end{equation}
This is a special case because $\sigma$ is rarely known in practice.

For battery life,
\[
n=81,\qquad \bar{x}=9.2\text{ hours},\qquad
\sigma=0.9\text{ hours},\qquad z^*=1.96.
\]
Step by step,
\[
\SE(\bar{x})=\frac{0.9}{\sqrt{81}}=0.1,
\qquad \mathrm{ME}=1.96(0.1)=0.196,
\]
so
\[
9.2\pm0.196=(9.004,9.396)\text{ hours}.
\]
We are $95\%$ confident that the population mean battery life lies between
$9.004$ and $9.396$ hours.  More precisely, if this sampling and interval
procedure were repeated many times, approximately $95\%$ of the constructed
intervals would contain the true $\mu$.

\subsection{A one-sample \texorpdfstring{$t$}{t}-interval for a mean}
When $\sigma$ is unknown and is estimated by $s$, use
\begin{equation}\label{eq:tinterval}
\bar{x}\pm t^*_{n-1}\frac{s}{\sqrt n},
\qquad df=n-1.
\end{equation}
For an illustrative sample of $n=16$ batteries with $\bar{x}=9.2$ hours and
$s=0.9$ hours, a $95\%$ interval uses $df=15$ and
$t^*_{15}\approx2.131$.  Thus
\[
\SE(\bar{x})=\frac{0.9}{\sqrt{16}}=0.225,
\qquad \mathrm{ME}=2.131(0.225)=0.4795\approx0.480,
\]
so the interval is
\[
9.2\pm0.480=(8.720,9.680)\text{ hours}.
\]
Under the design and model conditions, this method captures the true mean in
about $95\%$ of repeated samples.

\subsection{A one-proportion interval}
For a population proportion, an introductory large-sample interval is
\begin{equation}\label{eq:prop-interval}
\widehat p\pm z^*\sqrt{\frac{\widehat p(1-\widehat p)}{n}}.
\end{equation}
If $56$ of $100$ independent sampled students attended a review session, then
$\widehat p=0.56$.  The success--failure counts are $56$ and $44$, both at
least $10$.  A $95\%$ interval is
\[
0.56\pm1.96\sqrt{\frac{0.56(0.44)}{100}}
=0.56\pm1.96(0.04964)=0.56\pm0.0973,
\]
so the interval is approximately $(0.463,0.657)$.  We are $95\%$ confident
that the population proportion of students who would attend review lies
between $46.3\%$ and $65.7\%$, provided the sampling design supports that
population conclusion.

\subsection{Conditions, precision, and sample size}
For means, check that the data come from a random sample or randomized
process, observations are independent, and the sampling distribution of the
mean is approximately normal.  Approximate normality follows from a normal
population or from a sufficiently large sample under the central limit
theorem when extreme skew and outliers are not dominant.  When sampling
without replacement, a common independence check is $n\le10\%$ of the
population.  For proportions, check randomization, independence, a binary
outcome, and sufficient success--failure counts.

Higher confidence increases the critical value and interval width; larger
$n$ reduces standard error at rate $1/\sqrt n$.  To plan a margin of error
$M$ for a mean using a planning value of $\sigma$,
\begin{equation}\label{eq:n-mean}
n=\left(\frac{z^*\sigma}{M}\right)^2.
\end{equation}
For a proportion,
\begin{equation}\label{eq:n-prop}
n=\widehat p(1-\widehat p)\left(\frac{z^*}{M}\right)^2.
\end{equation}
If no prior estimate is available, $\widehat p=0.5$ is conservative.  Planned
sample sizes are rounded up.

\begin{keybox}
Confidence concerns the reliability of a method under repeated sampling, not
the percentage of observations in the interval.  A narrow interval can still
be misleading when sampling is biased or assumptions fail.  Confidence
intervals and hypothesis tests are complementary tools for inference.
\end{keybox}

\section{Hypothesis Testing}\label{sec:testing}
A \emph{hypothesis test} uses sample evidence to evaluate a claim about a
population.  The null hypothesis $H_0$ states the model being tested; the
alternative $H_a$ describes the departure of interest.  The same inferential
workflow applies: state the question, identify the parameter, check design
and conditions, compute a statistic, and interpret in context.

\subsection{Logic and terminology}
Choose a significance level $\alpha$, the probability of rejecting $H_0$ when
it is true (a Type I error).  A Type II error is failing to reject a false
$H_0$; power is
\begin{equation}\label{eq:power}
\text{power}=1-\beta.
\end{equation}
A test statistic measures discrepancy between data and $H_0$.  The $p$-value
is computed from the null distribution of the selected test statistic in the
direction or directions specified by $H_a$; it is the probability, assuming
$H_0$, of a result at least as inconsistent with $H_0$ as the observed
result.  Reject $H_0$ when $p\le\alpha$; equivalently, reject when the test
statistic falls in the rejection region determined by $\alpha$.  Otherwise
\emph{fail to reject} $H_0$.  A $p$-value is neither the probability that
$H_0$ is true nor an effect size.

Alternatives determine tail areas:
\[
H_a:\mu>\mu_0\quad\text{right-tailed},\qquad
H_a:\mu<\mu_0\quad\text{left-tailed},\qquad
H_a:\mu\ne\mu_0\quad\text{two-tailed}.
\]
Power generally increases with larger true effects, larger samples, lower
variability, or larger $\alpha$, but increasing $\alpha$ also increases Type
I error risk.  Statistical significance and practical importance differ: a
tiny effect can be significant in a very large sample, while an important
effect can fail to reach significance in a small or noisy sample.

\subsection{One-sample \texorpdfstring{$t$}{t}-test for a mean}
When $\sigma$ is unknown, a one-sample mean test uses
\begin{equation}\label{eq:t-test}
t=\frac{\bar{x}-\mu_0}{s/\sqrt n},
\qquad df=n-1.
\end{equation}
Using the illustrative battery sample from Section~\ref{sec:ci}, test at
$\alpha=0.05$ whether the mean battery life differs from $9$ hours:
\[
H_0:\mu=9,
\qquad H_a:\mu\ne9.
\]
With $n=16$, $\bar{x}=9.2$, and $s=0.9$,
\[
t=\frac{9.2-9}{0.9/\sqrt{16}}=\frac{0.2}{0.225}=0.889,
\qquad df=15.
\]
A two-tailed $t$ calculation gives $p\approx0.39$.  Since $p>0.05$, fail to
reject $H_0$.  The sample does not provide strong evidence that the population
mean battery life differs from $9$ hours; this does not prove the mean is
exactly $9$ hours.  Conditions are the same design, independence, and
shape/normality checks used for the one-sample $t$-interval.

\subsection{One-proportion \texorpdfstring{$z$}{z}-test}
For testing $H_0:p=p_0$, use the null-based standard error
\begin{equation}\label{eq:prop-test}
z=\frac{\widehat p-p_0}{\sqrt{p_0(1-p_0)/n}}.
\end{equation}
This differs from the confidence-interval standard error, which uses
$\widehat p$ because the interval estimates the unknown $p$.

For the review-attendance example, test
\[
H_0:p=0.50,
\qquad H_a:p\ne0.50,
\]
with $\widehat p=56/100=0.56$.  The null success--failure counts are
$n p_0=50$ and $n(1-p_0)=50$, both sufficient for the normal approximation.
Then
\[
z=\frac{0.56-0.50}{\sqrt{0.50(0.50)/100}}=\frac{0.06}{0.05}=1.20.
\]
The two-tailed $p$-value is approximately $0.23$.  At $\alpha=0.05$, fail to
reject $H_0$; the data do not provide strong evidence that the population
attendance proportion differs from $0.50$.

\subsection{Chi-square goodness-of-fit test}
For categorical counts, the chi-square statistic is
\begin{equation}\label{eq:chisq}
\chi^2=\sum_{i=1}^{k}\frac{(O_i-E_i)^2}{E_i},
\end{equation}
where $k$ is the number of categories, $O_i$ is observed frequency, $E_i$ is
the expected frequency under $H_0$, and $\chi^2$ is the total standardized
squared discrepancy.  If no parameters are estimated from the data, the
degrees of freedom are $df=k-1$.

For a hypothetical set of $60$ die rolls, test
\[
H_0:\ p_1=\cdots=p_6=\frac16
\quad\text{against}\quad
H_a:\ \text{at least one face probability differs}.
\]
Each expected count is $E_i=60(1/6)=10$.

\begin{table}[ht]
\centering
\caption{Reproducible chi-square goodness-of-fit calculation.}\label{tab:chisq-calc}
\begin{tabularx}{0.85\textwidth}{@{}lYYYY@{}}
\toprule
Face & Observed $O_i$ & Expected $E_i$ & Contribution $(O_i-E_i)^2/E_i$ & Pearson residual \\
\midrule
1 & 14 & 10 & $16/10=1.6$ & $4/\sqrt{10}=1.26$ \\
2 & 12 & 10 & $4/10=0.4$ & $2/\sqrt{10}=0.63$ \\
3 & 7 & 10 & $9/10=0.9$ & $-3/\sqrt{10}=-0.95$ \\
4 & 8 & 10 & $4/10=0.4$ & $-2/\sqrt{10}=-0.63$ \\
5 & 9 & 10 & $1/10=0.1$ & $-1/\sqrt{10}=-0.32$ \\
6 & 10 & 10 & $0/10=0.0$ & $0$ \\
\midrule
Total & 60 & 60 & $3.4$ & \\
\bottomrule
\end{tabularx}
\end{table}

With $k=6$, $df=5$.  At $\alpha=0.05$, the critical value is approximately
$11.07$, and the right-tail $p$-value for $\chi^2=3.4$ is approximately
$0.64$.  Since $3.4<11.07$ and $0.64>0.05$, fail to reject $H_0$.  There is
insufficient evidence that the die is unfair; this does not prove fairness.
Pearson residuals help identify which categories contribute most to the
global statistic.

\subsection{Conditions and practical meaning}
For one-sample mean tests, check randomization, independence, and a normal
population or a sufficiently large sample without dominant skew or outliers.
For one-proportion tests, check randomization, independence, binary outcomes,
and null-based success--failure counts.  For chi-square goodness-of-fit tests,
check frequency counts from random, independent observations in mutually
exclusive categories and expected counts sufficiently large; a common
introductory guideline is that all $E_i\ge5$.  If probabilities or parameters
are estimated from the sample, chi-square degrees of freedom must be adjusted.

\begin{keybox}[Common cautions]
State hypotheses before inspecting results, verify conditions, and interpret
the conclusion in context.  Never ``accept'' $H_0$: a nonsignificant result
may reflect limited information.  The next section shifts from categorical
fit to quantitative relationships.
\end{keybox}

\section{Regression and Correlation}\label{sec:regression}
\emph{Correlation} quantifies the direction and strength of a linear
association between two quantitative variables.  \emph{Regression} models the
conditional mean or predicted response together with unexplained variation;
it is not a deterministic function that exactly determines every outcome.

\subsection{Correlation and scatterplots}
For paired observations $(x_i,y_i)$, Pearson's sample correlation is
\begin{equation}\label{eq:r}
r=\frac{1}{n-1}\sum_{i=1}^{n}
\left(\frac{x_i-\bar{x}}{s_x}\right)
\left(\frac{y_i-\bar{y}}{s_y}\right).
\end{equation}
Here $n$ is the number of pairs; $\bar{x},\bar{y}$ and $s_x,s_y$ are the
sample means and standard deviations; and $r\in[-1,1]$.  Its sign gives the
direction, while $|r|$ measures the strength of a \emph{linear} association.
It is unitless, unchanged by shifting or positive rescaling, and sensitive to
outliers.  Correlation does not by itself establish causation: lurking
variables, reverse direction, and study design matter.

Use a scatterplot before computing or trusting $r$.  The following
hypothetical paired data will be used throughout this section.

\begin{table}[ht]
\centering
\caption{Hypothetical study-hours and exam-score data.}\label{tab:reg-data}
\begin{tabularx}{0.7\textwidth}{@{}lYYYYYYYY@{}}
\toprule
Study hours $x$ & 1 & 2 & 3 & 4 & 5 & 6 & 7 & 8 \\
Score $y$ & 55 & 58 & 65 & 67 & 72 & 74 & 81 & 83 \\
\bottomrule
\end{tabularx}
\end{table}

A compact scatterplot listing is
\[
(1,55),\ (2,58),\ (3,65),\ (4,67),\ (5,72),\ (6,74),\ (7,81),\ (8,83).
\]
It shows a positive, roughly linear association with no obvious isolated
outlier in this small illustrative data set.  Direct calculation gives
\[
\bar{x}=4.5,
\quad \bar{y}=69.375,
\quad s_x\approx2.449,
\quad s_y\approx10.070,
\quad r\approx0.9932.
\]

\subsection{Least-squares regression and residuals}
The simple linear regression model is
\begin{equation}\label{eq:regline}
\widehat y=b_0+b_1x,
\end{equation}
where $x$ is the explanatory value, $\widehat y$ is the predicted response,
$b_0$ is the predicted response at $x=0$ when that interpretation is
meaningful, and $b_1$ is the predicted change in $y$ for a one-unit increase
in $x$.  In simple linear regression with an intercept,
\begin{equation}\label{eq:slope-intercept}
b_1=r\frac{s_y}{s_x},
\qquad b_0=\bar y-b_1\bar x.
\end{equation}
For the paired data,
\[
b_1\approx0.9932\left(\frac{10.070}{2.449}\right)=4.083,
\qquad b_0=69.375-4.083(4.5)=51.000.
\]
Thus
\[
\widehat y=51.000+4.083x.
\]
At $x=4$ study hours, the predicted score is
\[
\widehat y=51.000+4.083(4)=67.333.
\]
A residual is
\begin{equation}\label{eq:residual}
e_i=y_i-\widehat y_i.
\end{equation}
For the $x=4$ observation, $e=67-67.333=-0.333$ points.  Least squares chooses
$b_0,b_1$ to minimize
\begin{equation}\label{eq:least-squares}
\sum_{i=1}^{n}e_i^2.
\end{equation}

\begin{table}[ht]
\centering
\caption{Fitted values and residuals for the regression example.}\label{tab:residuals}
\begin{tabularx}{0.85\textwidth}{@{}lYYYY@{}}
\toprule
Observation & $x_i$ & $y_i$ & $\widehat y_i$ & $e_i=y_i-\widehat y_i$ \\
\midrule
1 & 1 & 55 & 55.083 & $-0.083$ \\
2 & 2 & 58 & 59.167 & $-1.167$ \\
3 & 3 & 65 & 63.250 & $1.750$ \\
4 & 4 & 67 & 67.333 & $-0.333$ \\
5 & 5 & 72 & 71.417 & $0.583$ \\
6 & 6 & 74 & 75.500 & $-1.500$ \\
7 & 7 & 81 & 79.583 & $1.417$ \\
8 & 8 & 83 & 83.667 & $-0.667$ \\
\bottomrule
\end{tabularx}
\end{table}

\subsection{Residual diagnostics, influence, and \texorpdfstring{$R^2$}{R2}}
Residual plots help assess whether a line is appropriate.  Random scatter
around zero supports the linear form.  A curved residual pattern suggests
nonlinearity; changing vertical spread suggests nonconstant variance;
unusually large residuals suggest response outliers; a time-ordered pattern
may suggest dependence.

A response outlier is unusual in the $y$ direction.  A high-leverage
observation has an unusual $x$ value.  An influential observation is one whose
removal materially changes the fitted line; high leverage can make a point
influential, especially when its residual is also large.

In simple linear regression with an intercept,
\begin{equation}\label{eq:r-squared}
R^2=r^2.
\end{equation}
Here $R^2\approx(0.9932)^2\approx0.9865$: about $98.65\%$ of the observed
variation in exam score is accounted for by its fitted linear relationship
with study hours in this illustrative data set.  This does not mean $98.65\%$
of scores are exactly predicted or that study hours cause that share of
variation.

\begin{table}[ht]
\centering
\caption{Association, prediction, explanation, and causation.}\label{tab:association-causation}
\begin{tabularx}{\textwidth}{@{}lYY@{}}
\toprule
Use & Statistical question & Caution \\
\midrule
Descriptive association & Do two variables move together in these data? & Does not require a causal story \\
Prediction & What response is expected for a given $x$? & Extrapolation outside the data range is risky \\
Explanatory modeling & How is the response related to explanatory variables? & Confounding may remain \\
Causal inference & What would happen under intervention? & Requires appropriate design, often random assignment \\
\bottomrule
\end{tabularx}
\end{table}

\begin{keybox}[Regression checklist]
Inspect the scatterplot for form, direction, strength, clusters, and unusual
points.  Check residuals for random scatter around zero, roughly constant
variance, outliers, and dependence.  Avoid extrapolation.  For regression
inference, also check the relevant random sampling or experimental design and
any needed normal-error condition.  Causal conclusions require an appropriate
design; correlation or a fitted line alone is not enough.
\end{keybox}

\section*{Final Summary}\addcontentsline{toc}{section}{Final Summary}
Statistical work begins by identifying observational units, variables, data
types, and study design.  Sampling design determines what population
conclusions may generalize to, while random assignment, not correlation
alone, is central to causal conclusions.  Center, spread, and graphs describe
a sample without concealing shape, outliers, or context.  Probability turns
uncertain outcomes into models: discrete variables use probability masses and
continuous variables use areas under density curves.  Sampling distributions
and standard errors connect probability models to inference.  Confidence
intervals estimate population parameters or effects with uncertainty, and
hypothesis tests quantify evidence against a specified null model; statistical
significance and practical importance are different questions.  Regression
and correlation describe linear association and prediction while requiring
careful attention to scatterplots, residuals, influential observations,
design, and causal limitations.  Assumptions and diagnostic checks are part
of statistical procedures, not optional afterthoughts.

\section*{Compact Formula Reference}\addcontentsline{toc}{section}{Compact Formula Reference}
\begin{center}
\small
\begin{tabularx}{\textwidth}{@{}lYY@{}}
\toprule
Concept & Formula & Conditions/notes \\
\midrule
Population/sample mean & $\displaystyle \mu=\frac1N\sum_{i=1}^{N}x_i$; $\displaystyle \bar{x}=\frac1n\sum_{i=1}^{n}x_i$ & Quantitative data \\
Weighted/frequency mean & $\displaystyle \bar{x}_w=\frac{\sum_iw_ix_i}{\sum_iw_i}$; $\displaystyle \bar{x}=\frac{\sum_if_ix_i}{\sum_if_i}$ & Weights/frequencies nonnegative, not all zero \\
Range, IQR, fences & $x_{\max}-x_{\min}$; $\mathrm{IQR}=Q_3-Q_1$; $Q_1\pm1.5\mathrm{IQR}$ & State quartile and box-plot convention \\
Variance & $\displaystyle \sigma^2=\frac1N\sum (x_i-\mu)^2$; $\displaystyle s^2=\frac1{n-1}\sum (x_i-\bar{x})^2$ & $s^2$ estimates population variance under random sampling \\
Sample standardized value & $\displaystyle z_i=\frac{x_i-\bar{x}}{s}$ & Descriptive; requires $s>0$ \\
Empirical probability & $\displaystyle \widehat P(A)=\frac{\#A}{\#\text{ trials}}$ & Stable repeated trials for probability interpretation \\
Probability axioms & $0\le P(A)\le1$; $P(S)=1$; disjoint additivity & Foundation for probability rules \\
Conditional/multiplication & $\displaystyle P(A\mid B)=\frac{P(A\cap B)}{P(B)}$; $P(A\cap B)=P(A\mid B)P(B)$ & $P(B)>0$ \\
Independence & $P(A\cap B)=P(A)P(B)$ & Special case of multiplication \\
Total probability/Bayes & $P(A)=P(A\mid B)P(B)+P(A\mid B^c)P(B^c)$; $\displaystyle P(B\mid A)=\frac{P(A\mid B)P(B)}{P(A)}$ & $B,B^c$ partition; $P(A)>0$ \\
Binomial PMF & $\displaystyle P(X=k)=\binom nkp^k(1-p)^{n-k}$ & Fixed $n$, independent Bernoulli trials, constant $p$ \\
Poisson PMF & $\displaystyle P(X=k)=e^{-\lambda}\lambda^k/k!$ & Counts in specified interval; constant rate and independent disjoint intervals \\
Discrete CDF/mean/variance & $F(x)=\sum_{t\le x}p(t)$; $\E(X)=\sum_xxp(x)$; $\Var(X)=\sum_x(x-\mu_X)^2p(x)$ & Valid PMF sums to $1$ \\
Transformation rules & $\E(a+bX)=a+b\E(X)$; $\Var(a+bX)=b^2\Var(X)$; $\Var(X+Y)=\Var(X)+\Var(Y)$ & Last rule requires independence \\
PDF/CDF & $P(a\le X\le b)=\int_a^bf(x)\,dx$; $F(x)=\int_{-\infty}^xf(t)\,dt$ & $f\ge0$ and total area $1$ \\
Continuous mean/variance & $\E(X)=\int_{-\infty}^{\infty}xf(x)\,dx$; $\Var(X)=\int_{-\infty}^{\infty}(x-\mu)^2f(x)\,dx$ & When integrals exist \\
Standard normal score & $\displaystyle z=\frac{x-\mu}{\sigma}$ & Population/model standardization; $\sigma>0$ \\
Standard errors & $\SE(\bar{x})=\sigma/\sqrt n$ or $s/\sqrt n$; $\SE(\widehat p)=\sqrt{\widehat p(1-\widehat p)/n}$ & SD describes observations; SE describes statistics \\
Margin of error & $\mathrm{ME}=\text{critical value}\times\SE$ & Critical value depends on confidence level/procedure \\
Mean intervals & $\displaystyle \bar{x}\pm z^*\frac{\sigma}{\sqrt n}$; $\displaystyle \bar{x}\pm t^*_{n-1}\frac{s}{\sqrt n}$ & $z$: known $\sigma$; $t$: unknown $\sigma$, $df=n-1$ \\
Proportion interval & $\displaystyle \widehat p\pm z^*\sqrt{\frac{\widehat p(1-\widehat p)}{n}}$ & Random/independent binary data; success--failure condition \\
Sample size planning & $\displaystyle n=(z^*\sigma/M)^2$; $\displaystyle n=\widehat p(1-\widehat p)(z^*/M)^2$ & Round up; use $\widehat p=0.5$ if no prior estimate \\
Mean/proportion tests & $\displaystyle t=\frac{\bar{x}-\mu_0}{s/\sqrt n}$; $\displaystyle z=\frac{\widehat p-p_0}{\sqrt{p_0(1-p_0)/n}}$ & Mean: $df=n-1$; proportion test uses null-based SE \\
Chi-square GOF & $\displaystyle \chi^2=\sum_{i=1}^{k}\frac{(O_i-E_i)^2}{E_i}$; $df=k-1$ & $df=k-1$ only if no parameters estimated; expected counts large enough \\
Correlation & $\displaystyle r=\frac1{n-1}\sum_{i=1}^{n}\frac{x_i-\bar{x}}{s_x}\frac{y_i-\bar{y}}{s_y}$ & Paired quantitative data; $s_x,s_y>0$ \\
Regression & $\widehat y=b_0+b_1x$; $b_1=r(s_y/s_x)$; $b_0=\bar y-b_1\bar x$; $e_i=y_i-\widehat y_i$ & Simple linear regression with intercept \\
Least squares and $R^2$ & Minimize $\sum e_i^2$; $R^2=r^2$ & $R^2=r^2$ for simple linear regression with an intercept \\
\bottomrule
\end{tabularx}
\end{center}

\begin{keybox}[Responsible communication]
Report the data source, sampling or assignment method, procedures used,
software or tables used for critical values and $p$-values, assumptions,
uncertainty, and practical context.  Reproducible reporting makes statistical
claims easier to check.
\end{keybox}

\end{document}
